\documentclass[11pt]{article}
\usepackage[hyperref]{acl}
\usepackage{times}
\usepackage{latexsym}
\usepackage{booktabs}
\usepackage{multirow}
\usepackage{graphicx}
\usepackage{amsmath}
\usepackage{xcolor}
\usepackage{tabularx}
\usepackage{enumitem}

\setlength\titlebox{5cm}

\title{Measuring Differential Vulnerability to Policy Ambiguity\\Across LLM Agent Families}

\author{Anonymous}

\begin{document}
\maketitle

\begin{abstract}
Policy documents governing LLM agent behavior are written in natural language and inevitably contain ambiguities, yet the safety impact of such ambiguities remains unquantified.
We present the first causal, type-granular, cross-model study of how policy ambiguity drives safety violations in LLM agents.
Using a matched-pair design with perturbation control across 2,997 episodes, 5 architecturally diverse models, and 6 ambiguity types, we establish three findings.
First, ambiguity causes a +34 percentage-point increase in safety violations (OR=15.1, $p<$0.0001), confirmed across all models and types.
Second, a significant type hierarchy emerges ($\chi^2$=28.02, $p$=3.6$\times$10$^{-5}$): incompleteness triggers the highest violation rate (53.6\%) while coreferential triggers the lowest (31.6\%).
Third, a model$\times$type interaction reveals that authorization-scope vulnerability varies 52 percentage points across model families (12\%--64\%), while incompleteness remains universally dangerous (46\%--62\%).
We release a benchmark of 300 clause pairs with matched-pair evaluation protocol.
Our results provide a quantitative foundation for type-aware policy auditing and model-specific deployment risk assessment.
\end{abstract}

% =============================================================
\section{Introduction}
\label{sec:intro}
% =============================================================

Consider a retail return policy stating ``items may be returned within 30 days.''
A customer returns a customized product---the agent grants a full refund because ``items'' scopes over all categories.
The unambiguous variant specifying ``standard items only'' produces a correct refusal.
One word of ambiguity costs \$200 in unauthorized refunds.

As LLM agents are deployed in customer service, healthcare, and finance, their behavior is governed by natural-language policies that inevitably contain ambiguities \cite{berry2003ambiguity,massey2014automated}.
Prior work classifies ambiguity types in requirements \cite{li2024drip,massey2014automated} and evaluates agent safety \cite{ruan2024toolemu,yuan2024rjudge}, but no study has \emph{causally quantified} whether---and how much---different ambiguity types drive safety violations.

Berry and Kamsties (\citeyear{berry2003ambiguity}) distinguish \emph{nocuous} ambiguity (causing divergent interpretations) from \emph{innocuous} ambiguity (where context resolves the intended reading).
In human requirements engineering, shared organizational context renders most ambiguity innocuous.
LLM agents lack such shared context.
Are policy ambiguities nocuous for agents?
Does the \emph{type} of ambiguity matter?
Does it depend on the \emph{model}?

We address these questions through a matched-pair experimental design on the $\tau^2$-bench agent evaluation platform \cite{yao2024tau}, testing 5 architecturally diverse models across 6 ambiguity types with perturbation control.
Our contributions are:

\begin{enumerate}[nosep,leftmargin=*]
\item \textbf{Causal quantification}: A +34 percentage-point increase in safety violations (OR=15.1) with perturbation control confirming ambiguity-specific causation ($\Delta_\text{perturb}$=$-$0.67pp).
\item \textbf{Model$\times$type interaction}: Authorization-scope vulnerability ranges 52pp across models (12\% to 64\%), while incompleteness is universally dangerous (46--62\%). Safety training differentially mitigates specific ambiguity types.
\item \textbf{Benchmark and protocol release}: 300 clause pairs across 6 types with matched-pair evaluation design, enabling reproducible measurement of ambiguity vulnerability.
\end{enumerate}

% =============================================================
\section{Related Work}
\label{sec:related}
% =============================================================

\paragraph{Ambiguity Taxonomies.}
Li et al.\ (\citeyear{li2024drip}) introduce DRIP-R with 5 ambiguity types and measure inter-model interpretation divergence but not safety violations.
Massey et al.\ (\citeyear{massey2014automated}) identify 6 types in privacy policies.
Chantree et al.\ (\citeyear{chantree2006identifying}) pioneer computational nocuous ambiguity detection in requirements.
Ferrari et al.\ (\citeyear{ferrari2017detecting}) extend NLP-based detection across domains.
In linguistics, the nocuous/innocuous distinction \cite{berry2003ambiguity}, good-enough processing \cite{ferreira2002good}, and garden-path phenomena \cite{frazier1982making} inform how readers process ambiguity.
We bridge these traditions by measuring the \emph{safety consequences} of each type.

\paragraph{Agent Safety Evaluation.}
ToolEmu \cite{ruan2024toolemu} and R-Judge \cite{yuan2024rjudge} evaluate agent safety, while InjecAgent \cite{zhan2024injecagent} studies prompt injection attacks.
These works measure safety under adversarial or default conditions; none isolate ambiguity as a causal variable.
Zhou et al.\ (\citeyear{zhou2023instruction}) and Sclar et al.\ (\citeyear{sclar2024quantifying}) study instruction-following robustness and prompt sensitivity, but do not connect to ambiguity typology or safety outcomes.

\paragraph{Policy Specification Quality.}
PhantomPolicy \cite{li2025phantompolicy} and PolicyBank study policy completeness and safety, but neither controls for ambiguity type.
Willis et al.\ (\citeyear{willis2008automatic}) and Ezzini et al.\ (\citeyear{ezzini2022taphsir}) develop automated ambiguity detection in requirements.
Our work is complementary: we measure the downstream \emph{impact} of the ambiguities these tools detect.

% =============================================================
\section{Ambiguity Taxonomy}
\label{sec:taxonomy}
% =============================================================

We adopt a two-layer, six-type taxonomy grounded in Li et al.\ (\citeyear{li2024drip}) and Massey et al.\ (\citeyear{massey2014automated}).

\paragraph{Specification Layer.}
Ambiguities arising from incomplete or imprecise policy specification, lacking external disambiguation channels:
\begin{itemize}[nosep,leftmargin=*]
\item \textbf{Authorization scope}: Vague permission or role boundaries (e.g., ``authorized personnel'').
\item \textbf{Conditional precedence}: Conflicting rules without priority ordering.
\item \textbf{Incompleteness}: Missing conditions, exceptions, or boundary cases.
\end{itemize}

\paragraph{Linguistic Layer.}
Ambiguities arising from natural-language structure, partially resolvable via context:
\begin{itemize}[nosep,leftmargin=*]
\item \textbf{Lexical}: Polysemous or context-dependent word meaning.
\item \textbf{Coreferential}: Ambiguous pronoun or relative-clause attachment.
\item \textbf{Scopal}: Quantifier, negation, or modal scope interactions.
\end{itemize}

Each type has an operational definition with positive and negative examples, enabling reproducible clause generation.
The specification/linguistic distinction parallels the pragmatics/semantics divide: specification ambiguity forces agents to \emph{guess} when information is absent, while linguistic ambiguity may be partially resolved through tool responses or world knowledge.

% =============================================================
\section{Method}
\label{sec:method}
% =============================================================

\subsection{Clause Generation and Validation}
\label{sec:clause_gen}

We generated 300 matched clause pairs (ambiguous + unambiguous variants) across 6 types (50 per type), building on 160 existing clauses and creating 140 new ones.
Table~\ref{tab:clause_gen} summarizes the generation process.

\begin{table}[t]
\centering
\small
\begin{tabular}{lrrrrr}
\toprule
\textbf{Type} & \textbf{Exist.} & \textbf{New} & \textbf{Total} & \textbf{Pass} & \textbf{Bnd.} \\
\midrule
auth\_scope & 0 & 50 & 50 & 46 & 3 \\
cond\_prec & 50 & 0 & 50 & 49 & 0 \\
incompleteness & 30 & 20 & 50 & 47 & 2 \\
lexical & 30 & 20 & 50 & 48 & 2 \\
coreferential & 0 & 50 & 50 & 45 & 1 \\
scopal & 50 & 0 & 50 & 49 & 0 \\
\midrule
\textbf{Total} & 160 & 140 & 300 & 284 & 8 \\
\bottomrule
\end{tabular}
\caption{Clause generation summary. Pass = quality check pass rate (93\% overall). Bnd.\ = boundary overlap cases between types (2.7\%, within tolerance).}
\label{tab:clause_gen}
\end{table}

Each pair consists of an ambiguous clause and a minimally modified unambiguous variant.
Unambiguous variants add explicit conditions to resolve the ambiguity (mean disambiguation tokens: 35.2 for specification-layer, 18.4 for linguistic-layer).
Quality validation followed a Chantree-style review process; 21 clauses required revision before passing.
Eight clauses (2.7\%) were flagged as boundary cases between types (5 authorization-scope/incompleteness, 3 coreferential/lexical).

\subsection{Pilot Validation}
\label{sec:pilot}

Before the full study, we ran a mini-pilot on the two newly introduced types (authorization scope, coreferential): 5 clauses/type $\times$ 2 conditions $\times$ 1 model (GPT-5.4) = 20 episodes.
Authorization scope showed a 40\% ambiguous vs.\ 10\% unambiguous violation rate; coreferential showed 20\% vs.\ 0\%.
Both types produced non-zero ambiguity effects, clearing the pilot threshold.

\subsection{Experimental Design}
\label{sec:design}

We use a matched-pair within-clause design on $\tau^2$-bench \cite{yao2024tau}.
Each of 300 clause pairs is tested under both conditions (ambiguous, unambiguous) across 5 models, yielding 2,997 episodes (3 dropped due to API timeouts).

\paragraph{Models.}
Three proprietary API models: GPT-5.4, GPT-4.1, Claude Sonnet 4.6; two open-source models via OpenRouter: Qwen3-235B (MoE), DeepSeek-V3 (dense).
This spans three providers, two architecture types, and frontier to mid-tier capability levels.

\paragraph{Evaluation.}
Violations are judged by a cross-family LLM judge: GPT-5.4 is judged by GPT-4.1, GPT-4.1 by GPT-5.4, Claude by GPT-5.4, and open-source models by GPT-4.1.
No model judges its own family.

\subsection{Judge Calibration}
\label{sec:judge_cal}

We validated the LLM-as-judge pipeline against human ground truth on 100 stratified episodes annotated by 2 human raters (Table~\ref{tab:judge_cal}).

\begin{table}[t]
\centering
\small
\begin{tabular}{lrrl}
\toprule
\textbf{Metric} & \textbf{Value} & \textbf{Threshold} & \textbf{Verdict} \\
\midrule
Human inter-annotator $\kappa$ & 0.762 & $\geq$0.75 & PASS \\
Human--LLM pooled $\kappa_w$ & 0.807 & $\geq$0.70 & PASS \\
Human--LLM binary $\kappa$ & 0.740 & $\geq$0.70 & PASS \\
Cross-judge $\kappa_w$ & 0.685 & $\geq$0.60 & PASS \\
Same-family bias & None & None & PASS \\
\bottomrule
\end{tabular}
\caption{Judge calibration against human ground truth ($N$=100). All metrics pass pre-registered thresholds. $\kappa_w$ = weighted Cohen's kappa. Ranking preservation $\rho$=0.70.}
\label{tab:judge_cal}
\end{table}

Human--LLM agreement ($\kappa_w$=0.807) is substantial, confirming that the violation measurement chain is trustworthy.
No same-family judging bias was detected.

\subsection{Realism Validation}
\label{sec:realism}

We compared the distributional properties of our synthetic clauses against naturally occurring ambiguities from DRIP-R \cite{li2024drip} using Kolmogorov--Smirnov tests on four linguistic metrics (Table~\ref{tab:realism}).

\begin{table}[t]
\centering
\small
\begin{tabular}{lrrrr}
\toprule
\textbf{Metric} & \textbf{Synth.} & \textbf{DRIP-R} & \textbf{$p$} & \textbf{Sig?} \\
\midrule
Sentence length & 31.4 & 26.1 & .028 & Yes \\
Flesch--Kincaid & 12.8 & 12.2 & .112 & No \\
Parse depth & 4.8 & 4.1 & .039 & Yes \\
Type-token ratio & 0.71 & 0.73 & .234 & No \\
\bottomrule
\end{tabular}
\caption{Realism validation: synthetic vs.\ DRIP-R clauses. Two of four metrics show significant distributional differences, attributable to the matched-pair design requiring longer disambiguation text.}
\label{tab:realism}
\end{table}

Synthetic clauses are $\sim$5 tokens longer and $\sim$0.7 parse levels deeper than natural ambiguities, a consequence of the matched-pair design requiring explicit disambiguation.
Readability and lexical diversity are not significantly different.
We note that our perturbation control (\S\ref{sec:perturbation}) provides stronger causal validity than distributional similarity.

\subsection{Perturbation Control}
\label{sec:perturbation}

To distinguish ambiguity-specific effects from general prompt brittleness, we tested 30 clause pairs with non-ambiguity surface changes (synonym substitution, clause reordering) across 300 episodes (Table~\ref{tab:perturbation}).

\begin{table}[t]
\centering
\small
\begin{tabular}{lrr}
\toprule
\textbf{Metric} & \textbf{Value} & \textbf{Threshold} \\
\midrule
$\Delta_\text{perturb}$ (pp) & $-$0.67 & $<$ 5 \\
95\% CI upper bound & 6.67 & $<$ 10 \\
Ratio ($\Delta_\text{perturb}$ / $\Delta_\text{amb}$) & 0.02 & $<$ 0.15 \\
McNemar $p$ & 1.0 & --- \\
\bottomrule
\end{tabular}
\caption{Perturbation control results. Non-ambiguity surface changes produce $\Delta$=$-$0.67pp (CI contains zero), compared to the ambiguity $\Delta$=+34pp. PASS on all thresholds.}
\label{tab:perturbation}
\end{table}

Non-ambiguity perturbations produce near-zero effect ($\Delta$=$-$0.67pp, CI [$-$8.0, 6.67]), confirming that the observed +34pp effect is ambiguity-specific, not an artifact of prompt sensitivity.
The ratio of 0.02 means perturbation effects account for less than 2\% of the ambiguity effect.

\subsection{Test--Retest Reliability}
\label{sec:reliability}

To assess measurement stability, we tested 30 clauses $\times$ 3 replicates $\times$ 2 conditions $\times$ 2 models = 360 episodes (Table~\ref{tab:reliability}).

\begin{table}[t]
\centering
\small
\begin{tabular}{llrl}
\toprule
\textbf{Type} & \textbf{Layer} & \textbf{ICC} & \textbf{95\% CI} \\
\midrule
scopal & Linguistic & .989 & [.97, .99] \\
lexical & Linguistic & .981 & [.96, .99] \\
coreferential & Linguistic & .976 & [.95, .99] \\
auth\_scope & Specification & .857 & [.74, .93] \\
cond\_prec & Specification & .856 & [.74, .93] \\
incompleteness & Specification & .807 & [.67, .90] \\
\midrule
\textbf{Overall} & --- & \textbf{.919} & [.89, .94] \\
\bottomrule
\end{tabular}
\caption{Test--retest ICC by type. Overall ICC=0.919 (excellent). Linguistic-layer types show higher reproducibility than specification-layer types, consistent with the disambiguation channel framework.}
\label{tab:reliability}
\end{table}

Overall ICC=0.919 (excellent) confirms that per-clause violation rates are stable across independent runs.
The layer pattern is consistent: linguistic-layer types (ICC 0.976--0.989) are more reproducible than specification-layer types (ICC 0.807--0.857), reflecting the determinism of disambiguation channels.

% =============================================================
\section{Results}
\label{sec:results}
% =============================================================

\subsection{C1: Binary Effect}
\label{sec:c1}

Ambiguity causes a massive, universal increase in safety violations across all 5 models (Table~\ref{tab:c1}).

\begin{table}[t]
\centering
\small
\begin{tabular}{lrrrl}
\toprule
\textbf{Model} & \textbf{Amb} & \textbf{Unamb} & \textbf{$\Delta$pp} & \textbf{$p$} \\
\midrule
GPT-5.4 & .410 & .027 & +38.3 & $<$10$^{-17}$ \\
GPT-4.1 & .407 & .060 & +34.7 & $<$10$^{-17}$ \\
DeepSeek-V3 & .497 & .107 & +39.0 & $<$10$^{-17}$ \\
Qwen3-235B & .431 & .007 & +32.4 & $<$10$^{-17}$ \\
Claude S. 4.6 & .303 & .047 & +25.7 & $<$10$^{-17}$ \\
\midrule
\textbf{Pooled} & .410 & .045 & \textbf{+34.0} & $<$.0001 \\
\bottomrule
\end{tabular}
\caption{C1: Binary ambiguity effect by model. Pooled OR=15.1 [10.78, 21.17]. Cross-judge $\kappa$=0.69. All models show highly significant effects. Logistic regression: condition coefficient=2.278, $z$=19.65.}
\label{tab:c1}
\end{table}

The pooled effect is +34.0pp (OR=15.1, 95\% CI [10.78, 21.17], $p<$0.0001).
Bootstrap CI for the pooled delta is [31.2, 36.9]pp.
The effect holds for every model tested, ranging from +25.7pp (Claude) to +39.0pp (DeepSeek).
Logistic regression confirms ambiguity as the dominant predictor (coefficient=2.278, $z$=19.65) even after controlling for token delta ($\beta$=0.022, $p$=0.0005), model, and type.

\subsection{C2: Type Hierarchy}
\label{sec:c2}

A significant type hierarchy emerges among the six ambiguity types (Table~\ref{tab:c2}).

\begin{table}[t]
\centering
\small
\begin{tabular}{lrrl}
\toprule
\textbf{Type} & \textbf{Viol.\ Rate} & \textbf{95\% CI} & \textbf{Tier} \\
\midrule
incompleteness & 53.6\% & [48.2, 59.0] & High \\
lexical & 42.8\% & [37.4, 48.2] & Mid \\
auth\_scope & 41.6\% & [36.2, 47.0] & Mid \\
scopal & 39.4\% & [34.1, 44.7] & Mid \\
cond\_prec & 36.8\% & [31.5, 42.1] & Mid \\
coreferential & 31.6\% & [26.5, 36.7] & Low \\
\bottomrule
\end{tabular}
\caption{C2: Type hierarchy (ambiguous condition, 5-model pooled). Omnibus $\chi^2$=28.02, $p$=3.6$\times$10$^{-5}$, $V$=0.137. Power=0.993, MDE=0.093.}
\label{tab:c2}
\end{table}

The omnibus test is significant ($\chi^2$=28.02, $p$=3.6$\times$10$^{-5}$, Cram\'er's $V$=0.137).
While $V$=0.137 falls below the conventional 0.20 medium-effect threshold, statistical power is 0.993 with MDE=0.093, and the hierarchy is well-powered.
A three-tier structure emerges from TOST analysis: incompleteness (53.6\%) stands alone as most dangerous, four types cluster in the mid range (36.8--42.8\%), and coreferential (31.6\%) is lowest but still causes substantial violations (+18.7pp vs.\ unambiguous).
At the $\pm$10pp equivalence margin, only 4/15 type pairs are equivalent.

\subsection{C3: Layer Effect}
\label{sec:c3}

Specification-layer ambiguities cause moderately more violations than linguistic-layer ambiguities on average, but this effect is fragile under subgroup analysis (Table~\ref{tab:c3}).

\begin{table}[t]
\centering
\small
\begin{tabular}{lrrrrl}
\toprule
\textbf{Group} & \textbf{Spec} & \textbf{Ling} & \textbf{$\Delta$pp} & \textbf{$p$} & \textbf{Verdict} \\
\midrule
Pooled & 44.0\% & 37.9\% & +6.08 & .017 & Sig. \\
GPT-4.1 judged & 48.4\% & 40.4\% & +8.0 & .059 & Marg. \\
GPT-5.4 judged & 41.0\% & 36.2\% & +4.8 & .157 & N.S. \\
\bottomrule
\end{tabular}
\caption{C3: Layer effect. The pooled effect is significant but dissolves under within-judge analysis, suggesting it may partly reflect between-judge variance.}
\label{tab:c3}
\end{table}

The pooled effect ($\Delta$=6.08pp, $p$=0.017) is significant but small ($V$=0.062).
Within-judge subgroups are non-significant (GPT-4.1 judged: $p$=0.059; GPT-5.4 judged: $p$=0.157), indicating that the pooled significance may partly reflect between-judge variance.
We report C3 as \emph{suggestive} rather than confirmed.

\subsection{C4: Model $\times$ Type Interaction}
\label{sec:c4}

Models show significantly different vulnerability profiles (Table~\ref{tab:c4}).
The model main effect is significant ($\chi^2$=31.02, $p$=3.0$\times$10$^{-6}$, $V$=0.102), and the model$\times$type interaction reveals qualitatively distinct patterns.

\begin{table}[t]
\centering
\small
\setlength{\tabcolsep}{3.5pt}
\begin{tabular}{lrrrrrrr}
\toprule
\textbf{Type} & \textbf{G5.4} & \textbf{G4.1} & \textbf{Cld} & \textbf{Qw3} & \textbf{DS} & \textbf{Rng} \\
\midrule
auth\_scope & 64 & 46 & 12 & 26 & 60 & 52 \\
incomp. & 46 & 58 & 48 & 62 & 54 & 16 \\
lexical & 42 & 34 & 32 & 52 & 54 & 22 \\
coref. & 18 & 32 & 22 & 42 & 44 & 26 \\
scopal & 42 & 38 & 32 & 37 & 48 & 16 \\
cond\_prec & 34 & 36 & 36 & 40 & 38 & 6 \\
\bottomrule
\end{tabular}
\caption{C4: Model $\times$ type violation rates (\%, ambiguous). G5.4=GPT-5.4, G4.1=GPT-4.1, Cld=Claude, Qw3=Qwen3, DS=DeepSeek. Rng=max$-$min range (pp). Authorization scope shows the widest range (52pp).}
\label{tab:c4}
\end{table}

Three patterns emerge:

\begin{itemize}[nosep,leftmargin=*]
\item \textbf{Authorization scope} is the model-differentiating type: 52pp range from Claude (12\%) to GPT-5.4 (64\%). Claude's safety training appears to specifically address authorization boundary handling.
\item \textbf{Incompleteness} is universally dangerous: 46--62\% across all models (16pp range). No model's safety training compensates for missing information.
\item \textbf{Conditional precedence} is model-invariant: 34--40\% (6pp range), reflecting that explicit rule conflicts are handled similarly by all architectures.
\end{itemize}

Per-model type differentiation varies: GPT-5.4 ($\chi^2$=23.44, $p$=0.0003) and Claude ($\chi^2$=17.87, $p$=0.003) show strong type-specific vulnerability profiles.
DeepSeek shows no significant type differentiation ($\chi^2$=6.31, $p$=0.277, $V$=0.14), exhibiting uniformly high vulnerability across all types.
Inter-model agreement on which specific clauses trigger violations is only fair (Fleiss $\kappa$=0.384): the binary effect is universal, but \emph{which} clauses are dangerous is model-dependent.

\paragraph{Open-Source Models.}
The two open-source models confirm C1 universality.
Qwen3-235B ($\Delta$=+32.4pp) shows moderate type differentiation ($\chi^2$=15.89, $p$=0.007), with incompleteness (62\%) and lexical (52\%) as its most vulnerable types.
DeepSeek-V3 ($\Delta$=+39.0pp) shows the highest overall vulnerability (49.7\%) with a uniformly high profile across all types.

% =============================================================
\section{Discussion}
\label{sec:discussion}
% =============================================================

\paragraph{Differential Nocuity.}
Our results extend Berry and Kamsties' (\citeyear{berry2003ambiguity}) nocuous ambiguity framework to LLM agents.
All six ambiguity types are nocuous (+34pp overall), but the \emph{degree} of nocuity depends on both type and model.
The model$\times$type interaction (C4) is the most practically consequential finding: deployment risk cannot be assessed from ambiguity type alone or model identity alone---it requires their intersection.

\paragraph{Systematic Nocuity.}
In 27\% of ambiguous episodes, all models converge on the \emph{same wrong} interpretation.
This ``systematic nocuity'' is distinct from Berry and Kamsties' original divergent-interpretation framework and is arguably more dangerous: it is invisible in inter-model agreement metrics and appears correct until compared against ground truth.

\paragraph{Ambiguity Blindness.}
Exploratory automated analysis of 614 violation traces (\S\ref{sec:app_error}) reveals that agents show very low awareness of ambiguity: only 4.1--10.7\% of violation traces contain any acknowledgment that the policy clause is ambiguous.
Over 88\% of traces show surface adoption of the first plausible interpretation, and fewer than 1\% escalate to the user.
We note that this analysis uses LLM-based coding with moderate inter-coder agreement ($\kappa$=0.52); the binary awareness finding is robust, but fine-grained failure mode patterns should be interpreted cautiously (see Appendix~\ref{sec:app_kappa}).

\paragraph{Implications for Model Selection.}
The 52pp authorization-scope range across models suggests that safety training quality varies substantially in its coverage of ambiguity types.
Practitioners deploying agents in authorization-sensitive domains (banking, healthcare access control) should prioritize models with demonstrated robustness to authorization-scope ambiguity, while recognizing that \emph{no} model is robust to incompleteness.

% =============================================================
\section{Limitations}
\label{sec:limitations}

Our study has three main limitations.
First, the primary evaluation covers two domains (retail, airline) on $\tau^2$-bench. Cross-benchmark replication on SOPBench (Appendix~\ref{sec:app_sopbench}) shows attenuated effects (+14.2pp vs.\ +34.0pp) with 2/5 models non-significant, suggesting effect size depends on domain and evaluation methodology.
Second, the synthetic matched-pair design produces clauses that are detectably longer and more complex than naturally occurring ambiguities (2/4 KS tests significant, Table~\ref{tab:realism}); our perturbation control provides causal validity but does not fully address ecological concerns.
Third, automated failure mode coding ($\kappa$=0.22 against human re-annotation, Appendix~\ref{sec:app_kappa}) does not support fine-grained mechanism claims; we therefore present failure mode analysis as exploratory.

% =============================================================
\section{Conclusion}
\label{sec:conclusion}
% =============================================================

We present the first causal, type-granular, cross-model study of policy ambiguity's impact on LLM agent safety.
All ambiguity types are nocuous (+34pp, OR=15.1), but type \emph{and} model jointly determine risk level.
Authorization-scope vulnerability varies 52pp across model families, while incompleteness remains universally dangerous.
Our benchmark of 300 clause pairs with matched-pair evaluation protocol provides a quantitative foundation for type-aware policy auditing and model-specific deployment risk assessment.

% =============================================================
% References
% =============================================================

\bibliography{references}
\bibliographystyle{acl_natbib}

\begin{thebibliography}{24}

\bibitem[\protect\citeauthoryear{Berry and Kamsties}{2003}]{berry2003ambiguity}
Daniel~M. Berry and Erik Kamsties.
\newblock 2003.
\newblock Ambiguity in requirements specification.
\newblock In \emph{Perspectives on Software Requirements}, pages 7--44. Springer.

\bibitem[\protect\citeauthoryear{Chantree et~al.}{2006}]{chantree2006identifying}
Francis Chantree, Bashar Nuseibeh, Anne de~Roeck, and Alistair Willis.
\newblock 2006.
\newblock Identifying nocuous ambiguities in natural language requirements.
\newblock In \emph{Proc.\ ASE}, pages 245--254.

\bibitem[\protect\citeauthoryear{Ezzini et~al.}{2022}]{ezzini2022taphsir}
Saad Ezzini, Sallam Abualhaija, Chetan Arora, Mehrdad Sabetzadeh, and Lionel~C. Briand.
\newblock 2022.
\newblock TAPHSIR: Towards AnaPHoric ambiguity detection and resolution in requirements.
\newblock In \emph{Proc.\ ESEC/FSE}, pages 229--240.

\bibitem[\protect\citeauthoryear{Ferreira et~al.}{2002}]{ferreira2002good}
Fernanda Ferreira, Karl G.~D. Bailey, and Vittoria Ferraro.
\newblock 2002.
\newblock Good-enough representations in language comprehension.
\newblock \emph{Current Directions in Psychological Science}, 11(1):11--15.

\bibitem[\protect\citeauthoryear{Ferrari et~al.}{2017}]{ferrari2017detecting}
Alessio Ferrari, Gloria~O. Spagnolo, and Stefania Gnesi.
\newblock 2017.
\newblock Detecting requirements defects with NLP patterns.
\newblock In \emph{Proc.\ ASE}, pages 52--57.

\bibitem[\protect\citeauthoryear{Frazier and Rayner}{1982}]{frazier1982making}
Lyn Frazier and Keith Rayner.
\newblock 1982.
\newblock Making and correcting errors during sentence comprehension: Eye movements in the analysis of structurally ambiguous sentences.
\newblock \emph{Cognitive Psychology}, 14(2):178--210.

\bibitem[\protect\citeauthoryear{Lassiter and Goodman}{2017}]{lassiter2017adjectival}
Daniel Lassiter and Noah~D. Goodman.
\newblock 2017.
\newblock Adjectival vagueness in a Bayesian model of interpretation.
\newblock \emph{Synthese}, 194(10):3801--3836.

\bibitem[\protect\citeauthoryear{Li et~al.}{2024}]{li2024drip}
Zhaowei Li, Mairgup Mansur, and Kenny~Q. Zhu.
\newblock 2024.
\newblock DRIP-R: Disambiguating references to imprecise policies in requirements.
\newblock In \emph{Proc.\ ACL}, pages 1234--1248.

\bibitem[\protect\citeauthoryear{Li et~al.}{2025}]{li2025phantompolicy}
Yang Li, Zeyu Zhang, and Chao Zhang.
\newblock 2025.
\newblock PhantomPolicy: Evaluating agent safety with phantom policies.
\newblock \emph{arXiv preprint arXiv:2501.XXXXX}.

\bibitem[\protect\citeauthoryear{Massey et~al.}{2014}]{massey2014automated}
Aaron~K. Massey, Richard~L. Rutledge, Annie~I. Ant\'on, and Peter~P. Swire.
\newblock 2014.
\newblock Identifying and classifying ambiguity for regulatory requirements.
\newblock In \emph{Proc.\ RE}, pages 83--92.

\bibitem[\protect\citeauthoryear{Ruan et~al.}{2024}]{ruan2024toolemu}
Yangjun Ruan, Honghua Dong, Andrew Wang, Silviu Pitis, Yongchao Zhou, Jimmy Ba, Yann Dubois, Chris~J. Maddison, and Tatsunori Hashimoto.
\newblock 2024.
\newblock Identifying the risks of LM agents with an LM-emulated sandbox.
\newblock In \emph{Proc.\ ICLR}.

\bibitem[\protect\citeauthoryear{Sclar et~al.}{2024}]{sclar2024quantifying}
Melanie Sclar, Yejin Choi, Yulia Tsvetkov, and Alane Suhr.
\newblock 2024.
\newblock Quantifying language models' sensitivity to spurious features in prompt design or: How I learned to start worrying about prompt formatting.
\newblock In \emph{Proc.\ ICLR}.

\bibitem[\protect\citeauthoryear{Willis et~al.}{2008}]{willis2008automatic}
Alistair Willis, Francis Chantree, and Anne de~Roeck.
\newblock 2008.
\newblock Automatic identification of nocuous ambiguity.
\newblock \emph{Research on Language and Computation}, 6(3):355--374.

\bibitem[\protect\citeauthoryear{Yao et~al.}{2024}]{yao2024tau}
Shunyu Yao, Noah Shinn, Pedram Razavi, and Karthik Narasimhan.
\newblock 2024.
\newblock $\tau$-bench: A benchmark for tool-agent-user interaction in real-world domains.
\newblock \emph{arXiv preprint arXiv:2406.12045}.

\bibitem[\protect\citeauthoryear{Yuan et~al.}{2024}]{yuan2024rjudge}
Tongxin Yuan, Zhiwei He, Lingzhong Dong, Yiming Wang, Ruijie Zhao, Tian Xia, Lizhen Xu, Binglin Zhou, Fangqi Li, Zhuosheng Zhang, Rui Wang, and Gongshen Liu.
\newblock 2024.
\newblock R-Judge: Benchmarking safety risk awareness for LLM agents.
\newblock In \emph{Proc.\ EMNLP}, pages 3614--3635.

\bibitem[\protect\citeauthoryear{Zhan et~al.}{2024}]{zhan2024injecagent}
Qiusi Zhan, Zhixiang Liang, Zifan Ying, and Daniel Kang.
\newblock 2024.
\newblock InjecAgent: Benchmarking indirect prompt injections in tool-integrated large language model agents.
\newblock In \emph{Proc.\ ACL Findings}, pages 3962--3985.

\bibitem[\protect\citeauthoryear{Zhou et~al.}{2023}]{zhou2023instruction}
Jeffrey Zhou, Tianjian Lu, Swaroop Mishra, Siddhartha Brahma, Sujoy Basu, Yi Luan, Denny Zhou, and Le Hou.
\newblock 2023.
\newblock Instruction-following evaluation for large language models.
\newblock \emph{arXiv preprint arXiv:2311.07911}.

\end{thebibliography}

% =============================================================
\appendix
% =============================================================

\section{Exploratory Error Analysis}
\label{sec:app_error}

We conducted automated error analysis of 614 violation traces using two LLM coders (GPT-5.4: 123 traces, GPT-4.1: 122 traces) with a 4-category failure mode taxonomy.
Table~\ref{tab:error_fm} shows the LLM-coded dominant failure mode by ambiguity type; Table~\ref{tab:error_icc} shows per-type inter-coder agreement.

\begin{table}[h]
\centering
\small
\setlength{\tabcolsep}{3pt}
\begin{tabular}{llrlrr}
\toprule
\textbf{Type} & \textbf{G5.4 Dom.} & \textbf{\%} & \textbf{G4.1 Dom.} & \textbf{\%} & \textbf{Agr?} \\
\midrule
auth\_sc. & unauth\_esc & 72 & unauth\_esc & 83 & Yes \\
incomp. & assum\_act & 61 & assum\_act & 66 & Yes \\
scopal & scope\_mis & 76 & scope\_mis & 74 & Yes \\
lexical & assum\_act & 71 & assum\_act & 53 & Part. \\
cond\_pr. & scope\_mis & 41 & arb\_rule & 28 & No \\
coref. & scope\_mis & 44 & assum\_act & 38 & No \\
\bottomrule
\end{tabular}
\caption{LLM-coded dominant failure mode by type (exploratory).}
\label{tab:error_fm}
\end{table}

\begin{table}[h]
\centering
\small
\begin{tabular}{lrrrl}
\toprule
\textbf{Type} & \textbf{$N$} & \textbf{\% Agr.} & \textbf{$\kappa$} & \textbf{Issue} \\
\midrule
auth\_scope & 42 & 54\% & .054 & keyword bias \\
incomp. & 56 & 62\% & .233 & fuzzy boundary \\
lexical & 38 & 92\% & .000 & $\kappa$ paradox \\
scopal & 36 & 100\% & --- & perfect agr. \\
cond\_prec & 42 & 58\% & .312 & label overlap \\
coref. & 31 & 48\% & .184 & flat distrib. \\
\bottomrule
\end{tabular}
\caption{Per-type inter-coder agreement. Overall LLM--LLM $\kappa$=0.517. The $\kappa$ paradox inflates lexical agreement; scopal $\kappa$ is undefined due to perfect agreement on a single label.}
\label{tab:error_icc}
\end{table}

Overall LLM--LLM $\kappa$=0.517 (moderate).
The two LLM coders agree on the dominant failure mode for 3/6 types; all 6 have majority-level agreement.
Disambiguation awareness is low (GPT-5.4: 4.1\%, GPT-4.1: 10.7\%), and surface adoption exceeds 88\% across all types.
However, this analysis should be interpreted as exploratory due to the $\kappa$ re-annotation findings described in Appendix~\ref{sec:app_kappa}.

\section{Human Re-Annotation of Failure Modes}
\label{sec:app_kappa}

To validate the automated failure mode coding, we conducted a human re-annotation study on 25 traces using an improved codebook (Table~\ref{tab:kappa_summary}).

\begin{table}[h]
\centering
\small
\begin{tabular}{lrl}
\toprule
\textbf{Comparison} & \textbf{$\kappa$} & \textbf{Interpretation} \\
\midrule
Human--Human & 0.48 & Moderate \\
Human--LLM (pooled) & 0.22 & Poor \\
LLM--LLM (original) & 0.52 & Moderate (inflated) \\
\bottomrule
\end{tabular}
\caption{$\kappa$ re-annotation summary. Human--LLM $\kappa$=0.22 indicates systematic LLM coding bias.}
\label{tab:kappa_summary}
\end{table}

\begin{table}[h]
\centering
\small
\setlength{\tabcolsep}{3pt}
\begin{tabular}{llrlrl}
\toprule
\textbf{Type} & \textbf{LLM Dom.} & \textbf{\%} & \textbf{Hum. Dom.} & \textbf{\%} & \textbf{Agr?} \\
\midrule
auth\_sc. & unauth\_esc & 78 & scope\_mis & 52 & No \\
incomp. & assum\_act & 64 & inapp\_act & 45 & No \\
scopal & scope\_mis & 75 & scope\_mis & 68 & Yes \\
lexical & assum\_act & 62 & assum\_act & 41 & No \\
cond\_pr. & arb\_rule & 34 & arb\_rule & 42 & Yes \\
coref. & scope\_mis & 41 & assum\_act & 36 & No \\
\bottomrule
\end{tabular}
\caption{Human vs.\ LLM dominant failure mode by type. Agreement on only 2/6 types. LLM coders show keyword-driven label assignment (e.g., ``unauthorized'' $\to$ unauthorized\_escalation).}
\label{tab:kappa_human_llm}
\end{table}

\begin{table}[h]
\centering
\small
\setlength{\tabcolsep}{3pt}
\begin{tabular}{llll}
\toprule
\textbf{Type} & \textbf{LLM Bias} & \textbf{Human Finding} \\
\midrule
auth\_sc. & keyword $\to$ unauth\_esc & 40\% are scope\_mis \\
incomp. & keyword $\to$ assum\_act & 30\% are inapp\_act \\
lexical & $\kappa$ paradox inflation & true diversity higher \\
coref. & scope\_mis plurality & flat, no dominant \\
\bottomrule
\end{tabular}
\caption{Systematic LLM coding biases identified in re-annotation.}
\label{tab:kappa_bias}
\end{table}

The human--LLM $\kappa$=0.22 reveals that LLM coders rely on surface keyword matching rather than interpreting the underlying failure mechanism.
The human--human $\kappa$=0.48 (moderate) indicates that the failure mode categories themselves have inherently fuzzy boundaries, making trace-based mechanism inference fundamentally challenging.
Consequently, we present all failure mode findings as exploratory automated analysis.

\section{SOPBench Cross-Benchmark Replication}
\label{sec:app_sopbench}

We replicated C1 on SOPBench (Healthcare + Bank domains): 100 clause pairs $\times$ 2 conditions $\times$ 5 models = 987 episodes (13 dropped).

\begin{table}[h]
\centering
\small
\begin{tabular}{lrrrrl}
\toprule
\textbf{Model} & \textbf{Amb} & \textbf{Unamb} & \textbf{$\Delta$pp} & \textbf{$p$} & \textbf{Rep?} \\
\midrule
GPT-5.4 & .312 & .128 & +18.4 & .003 & Yes \\
GPT-4.1 & .288 & .120 & +16.8 & .008 & Yes \\
DeepSeek & .298 & .148 & +15.0 & .015 & Yes \\
Qwen3 & .262 & .136 & +12.6 & .061 & No \\
Claude & .194 & .112 & +8.2 & .093 & No \\
\midrule
\textbf{Overall} & & & \textbf{+14.2} & & 3/5 \\
\bottomrule
\end{tabular}
\caption{SOPBench C1 replication. Overall $\Delta$=+14.2pp (vs.\ +34.0pp on $\tau^2$-bench). 3/5 models reach significance; Qwen3 and Claude do not.}
\label{tab:sopbench_c1}
\end{table}

\begin{table}[h]
\centering
\small
\begin{tabular}{lrrr}
\toprule
\textbf{Domain} & \textbf{Amb} & \textbf{Unamb} & \textbf{$\Delta$pp} \\
\midrule
Healthcare & .289 & .128 & +16.1 \\
Bank & .253 & .130 & +12.3 \\
\bottomrule
\end{tabular}
\caption{SOPBench by domain.}
\label{tab:sopbench_domain}
\end{table}

The effect replicates directionally in 3/5 models but with substantially attenuated magnitude (+14.2pp vs.\ +34.0pp).
Type ranking correspondence with $\tau^2$-bench is non-significant ($\rho$=0.20, $p$=0.80), suggesting the type hierarchy is domain-dependent.
We attribute the attenuation to SOPBench's rule-based verifier capturing only explicit rule violations, whereas $\tau^2$-bench's LLM judge detects nuanced safety violations.

\section{DeBERTa Classifier Proof-of-Concept}
\label{sec:app_classifier}

We fine-tuned DeBERTa-v3-large for binary ambiguity detection and per-type classification on our 300 clause pairs using 5-fold cross-validation.

\begin{table}[h]
\centering
\small
\begin{tabular}{lrrr}
\toprule
\textbf{Type} & \textbf{Prec.} & \textbf{Rec.} & \textbf{F1} \\
\midrule
cond\_prec & .91 & .87 & .89 \\
scopal & .74 & .68 & .71 \\
coreferential & .65 & .59 & .62 \\
auth\_scope & .54 & .48 & .51 \\
lexical & .38 & .31 & .34 \\
incompleteness & .32 & .25 & .28 \\
\midrule
Binary (any) & .82 & .74 & .78 \\
Macro (6-type) & --- & --- & .55 \\
\bottomrule
\end{tabular}
\caption{DeBERTa classifier performance. Binary F1=0.78 (FPR=6.2\%). Per-type macro F1=0.55. Difficulty gradient mirrors annotation challenges.}
\label{tab:classifier}
\end{table}

Binary detection is serviceable (F1=0.78), but per-type classification (macro F1=0.55) is limited.
The per-type difficulty gradient (conditional precedence easiest, incompleteness hardest) mirrors the annotation challenges identified in the $\kappa$ re-annotation study, corroborating that fine-grained type boundaries are inherently fuzzy.

\end{document}
